% --- Archon physics formalization source begin ---
% archon:physics
% archon:covers PhyXMiniProblems/problem_phyx_mini_0222.lean
% archon:source-report reports/phyx_mini/problem_phyx_mini_0222.source.json

\chapter{Physics problem phyx\_mini\_0222}
\label{ch:PhyXMiniProblems_problem_phyx_mini_0222}

\paragraph{Problem source.}
\#\# Physical scenario

Two loudspeakers are at opposite ends of a railroad car as it moves past a stationary observer at \textbackslash{}( 12.0\textbackslash{},\textbackslash{}mathrm\{m/s\} \textbackslash{}), as shown in figure.The speakers have identical sound frequencies of \textbackslash{}( 348\textbackslash{},\textbackslash{}mathrm\{Hz\} \textbackslash{}).He is between the speakers, at B.

\#\# Question

What is the beat frequency heard by the observer?

\#\# Answer choices

- A: A: 30Hz

- B: B: 18Hz

- C: C: 12Hz

- D: D: 24Hz

\#\# Figure caption (auxiliary; use the image as primary evidence)

The image depicts a scene with three individuals labeled as A, B, and C standing next to a set of tracks. There is a platform on rails carrying two large speakers and a person identified as B standing on it. The platform is moving to the right with a velocity indicated by an arrow labeled \textbackslash{}( v = 12.0 \textbackslash{}, \textbackslash{}text\{m/s\} \textbackslash{}).

- Person A is standing on the right side of the moving platform, observing it.
- Person C is on the left side, also observing the platform, and positioned such that the platform is moving away from them.
- Person B is standing on the platform between the two speakers.

The ground beneath the tracks is visible, with details showing different layers of soil and rock.

\#\# Dataset metadata

Wave Properties; Physical Model Grounding Reasoning, Multi-Formula Reasoning

\paragraph{Recorded answer/context.}
D: D: 24Hz

\paragraph{Figure/image path.}
/root/proposal\_for\_physic/science-mango/phyx\_mini\_run/phyx\_data/test\_image/222.png

\paragraph{Formalization target.}
create a compiling Lean file with sorry bodies at `PhyXMiniProblems/problem\_phyx\_mini\_0222.lean`. The Lean declarations must preserve the physical quantities, dimensions or dimensional roles, figure labels, governing-law hypotheses, and final relation expressed by this problem.
Use Mathlib/Physlib names found through LeanExplore where available. If a physics API is missing, introduce faithful local abstractions rather than scalar placeholder aliases.

\begin{theorem}[Physics formalization target]
\label{thm:physics:phyx_mini_0222:target}
\lean{PhyXMiniProblems.ProblemPhyXMini0222.beatFrequencyAtB_matches_recordedAnswerD}
\uses{def:physics:phyx-mini-0222:phyxminiproblems-problemphyxmini0222-frequencyinhertz, def:physics:phyx-mini-0222:phyxminiproblems-problemphyxmini0222-movingrailroadcarspeakersetup, def:physics:phyx-mini-0222:phyxminiproblems-problemphyxmini0222-matchesproblemandprimaryfigure, def:physics:phyx-mini-0222:phyxminiproblems-problemphyxmini0222-usesstandardstillair, def:physics:phyx-mini-0222:phyxminiproblems-problemphyxmini0222-hasphysicaldopplerparameters, def:physics:phyx-mini-0222:phyxminiproblems-problemphyxmini0222-satisfiesdirectionalmovingsourcedopplerlaws, def:physics:phyx-mini-0222:phyxminiproblems-problemphyxmini0222-satisfiesbeatfrequencylaw, def:physics:phyx-mini-0222:phyxminiproblems-problemphyxmini0222-recordedanswerchoice, def:physics:phyx-mini-0222:phyxminiproblems-problemphyxmini0222-matchesanswerchoice}
The assigned autoformalize agent should translate this physics problem into Lean declarations in the covered file, with theorem and lemma proof bodies written as `by sorry`.
\end{theorem}
\begin{proof}
This is an autoformalization task, not a proof task. Produce statements that can later be proved by the physics prover without weakening the physical model.
\end{proof}

% --- Archon declaration topology begin ---
\section{Declaration topology}
\begin{definition}[Acoustic Frequency]
  \label{def:physics:phyx-mini-0222:phyxminiproblems-problemphyxmini0222-acousticfrequency}
  \lean{PhyXMiniProblems.ProblemPhyXMini0222.AcousticFrequency}
  A nonnegative physical acoustic frequency, with inverse-time dimension
\end{definition}
\begin{proof}
  This is the declared model definition of Acoustic Frequency.
\end{proof}
\begin{definition}[Axial Position]
  \label{def:physics:phyx-mini-0222:phyxminiproblems-problemphyxmini0222-axialposition}
  \lean{PhyXMiniProblems.ProblemPhyXMini0222.AxialPosition}
  A signed physical coordinate on the horizontal axis of the figure
\end{definition}
\begin{proof}
  This is the declared model definition of Axial Position.
\end{proof}
\begin{definition}[Axial Velocity]
  \label{def:physics:phyx-mini-0222:phyxminiproblems-problemphyxmini0222-axialvelocity}
  \lean{PhyXMiniProblems.ProblemPhyXMini0222.AxialVelocity}
  A signed physical velocity along the figure's horizontal axis. Positive velocity points to the right
\end{definition}
\begin{proof}
  This is the declared model definition of Axial Velocity.
\end{proof}
\begin{definition}[frequency In Hertz]
  \label{def:physics:phyx-mini-0222:phyxminiproblems-problemphyxmini0222-frequencyinhertz}
  \lean{PhyXMiniProblems.ProblemPhyXMini0222.frequencyInHertz}
  \uses{def:physics:phyx-mini-0222:phyxminiproblems-problemphyxmini0222-acousticfrequency}
  Read a physical acoustic frequency in hertz
\end{definition}
\begin{proof}
  This is the declared model definition of frequency In Hertz.
\end{proof}
\begin{definition}[position In Meters]
  \label{def:physics:phyx-mini-0222:phyxminiproblems-problemphyxmini0222-positioninmeters}
  \lean{PhyXMiniProblems.ProblemPhyXMini0222.positionInMeters}
  \uses{def:physics:phyx-mini-0222:phyxminiproblems-problemphyxmini0222-axialposition}
  Read a signed axial position in metres
\end{definition}
\begin{proof}
  This is the declared model definition of position In Meters.
\end{proof}
\begin{definition}[axial Velocity In Meters Per Second]
  \label{def:physics:phyx-mini-0222:phyxminiproblems-problemphyxmini0222-axialvelocityinmeterspersecond}
  \lean{PhyXMiniProblems.ProblemPhyXMini0222.axialVelocityInMetersPerSecond}
  \uses{def:physics:phyx-mini-0222:phyxminiproblems-problemphyxmini0222-axialvelocity}
  Read a signed axial velocity in metres per second
\end{definition}
\begin{proof}
  This is the declared model definition of axial Velocity In Meters Per Second.
\end{proof}
\begin{definition}[speed In Meters Per Second]
  \label{def:physics:phyx-mini-0222:phyxminiproblems-problemphyxmini0222-speedinmeterspersecond}
  \lean{PhyXMiniProblems.ProblemPhyXMini0222.speedInMetersPerSecond}
  Read Physlib's nonnegative physical speed in metres per second
\end{definition}
\begin{proof}
  This is the declared model definition of speed In Meters Per Second.
\end{proof}
\begin{definition}[Figure Person]
  \label{def:physics:phyx-mini-0222:phyxminiproblems-problemphyxmini0222-figureperson}
  \lean{PhyXMiniProblems.ProblemPhyXMini0222.FigurePerson}
  The three people labelled in the primary figure
\end{definition}
\begin{proof}
  This is the declared model definition of Figure Person.
\end{proof}
\begin{definition}[Speaker Label]
  \label{def:physics:phyx-mini-0222:phyxminiproblems-problemphyxmini0222-speakerlabel}
  \lean{PhyXMiniProblems.ProblemPhyXMini0222.SpeakerLabel}
  The two loudspeakers attached to the railroad car
\end{definition}
\begin{proof}
  This is the declared model definition of Speaker Label.
\end{proof}
\begin{definition}[Railroad Car End]
  \label{def:physics:phyx-mini-0222:phyxminiproblems-problemphyxmini0222-railroadcarend}
  \lean{PhyXMiniProblems.ProblemPhyXMini0222.RailroadCarEnd}
  The ends of the railroad car as oriented in the primary figure
\end{definition}
\begin{proof}
  This is the declared model definition of Railroad Car End.
\end{proof}
\begin{definition}[Vehicle Label]
  \label{def:physics:phyx-mini-0222:phyxminiproblems-problemphyxmini0222-vehiclelabel}
  \lean{PhyXMiniProblems.ProblemPhyXMini0222.VehicleLabel}
  The carrier on which both sound sources are mounted
\end{definition}
\begin{proof}
  This is the declared model definition of Vehicle Label.
\end{proof}
\begin{definition}[Axial Direction]
  \label{def:physics:phyx-mini-0222:phyxminiproblems-problemphyxmini0222-axialdirection}
  \lean{PhyXMiniProblems.ProblemPhyXMini0222.AxialDirection}
  The two orientations along the one-dimensional figure axis
\end{definition}
\begin{proof}
  This is the declared model definition of Axial Direction.
\end{proof}
\begin{definition}[Moving Railroad Car Speaker Setup]
  \label{def:physics:phyx-mini-0222:phyxminiproblems-problemphyxmini0222-movingrailroadcarspeakersetup}
  \lean{PhyXMiniProblems.ProblemPhyXMini0222.MovingRailroadCarSpeakerSetup}
  \uses{def:physics:phyx-mini-0222:phyxminiproblems-problemphyxmini0222-acousticfrequency, def:physics:phyx-mini-0222:phyxminiproblems-problemphyxmini0222-axialposition, def:physics:phyx-mini-0222:phyxminiproblems-problemphyxmini0222-axialvelocity, def:physics:phyx-mini-0222:phyxminiproblems-problemphyxmini0222-figureperson, def:physics:phyx-mini-0222:phyxminiproblems-problemphyxmini0222-speakerlabel, def:physics:phyx-mini-0222:phyxminiproblems-problemphyxmini0222-railroadcarend, def:physics:phyx-mini-0222:phyxminiproblems-problemphyxmini0222-vehiclelabel, def:physics:phyx-mini-0222:phyxminiproblems-problemphyxmini0222-axialdirection}
  Independent physical quantities for the two-source Doppler experiment. The two received frequencies and beatFrequencyAtB are unknown physical frequencies. In particular, this setup does not assign the requested beat frequency a numerical value
\end{definition}
\begin{proof}
  This is the declared model definition of Moving Railroad Car Speaker Setup.
\end{proof}
\begin{definition}[speaker Velocity Relative To Air In Meters Per Second]
  \label{def:physics:phyx-mini-0222:phyxminiproblems-problemphyxmini0222-speakervelocityrelativetoairinmeterspersecond}
  \lean{PhyXMiniProblems.ProblemPhyXMini0222.speakerVelocityRelativeToAirInMetersPerSecond}
  \uses{def:physics:phyx-mini-0222:phyxminiproblems-problemphyxmini0222-axialvelocityinmeterspersecond, def:physics:phyx-mini-0222:phyxminiproblems-problemphyxmini0222-speakerlabel, def:physics:phyx-mini-0222:phyxminiproblems-problemphyxmini0222-movingrailroadcarspeakersetup}
  A speaker's signed velocity relative to the air, in metres per second
\end{definition}
\begin{proof}
  This is the declared model definition of speaker Velocity Relative To Air In Meters Per Second.
\end{proof}
\begin{definition}[observer BVelocity Relative To Air In Meters Per Second]
  \label{def:physics:phyx-mini-0222:phyxminiproblems-problemphyxmini0222-observerbvelocityrelativetoairinmeterspersecond}
  \lean{PhyXMiniProblems.ProblemPhyXMini0222.observerBVelocityRelativeToAirInMetersPerSecond}
  \uses{def:physics:phyx-mini-0222:phyxminiproblems-problemphyxmini0222-axialvelocityinmeterspersecond, def:physics:phyx-mini-0222:phyxminiproblems-problemphyxmini0222-movingrailroadcarspeakersetup}
  Observer B's signed velocity relative to the air, in metres per second
\end{definition}
\begin{proof}
  This is the declared model definition of observer BVelocity Relative To Air In Meters Per Second.
\end{proof}
\begin{definition}[Matches Problem And Primary Figure]
  \label{def:physics:phyx-mini-0222:phyxminiproblems-problemphyxmini0222-matchesproblemandprimaryfigure}
  \lean{PhyXMiniProblems.ProblemPhyXMini0222.MatchesProblemAndPrimaryFigure}
  \uses{def:physics:phyx-mini-0222:phyxminiproblems-problemphyxmini0222-frequencyinhertz, def:physics:phyx-mini-0222:phyxminiproblems-problemphyxmini0222-positioninmeters, def:physics:phyx-mini-0222:phyxminiproblems-problemphyxmini0222-axialvelocityinmeterspersecond, def:physics:phyx-mini-0222:phyxminiproblems-problemphyxmini0222-movingrailroadcarspeakersetup}
  Problem-statement and primary-figure data. The picture orders C, the rear speaker, B, the front speaker, and A from left to right. Both speakers are fixed to the car, have the same 348 Hz emission, and move right with the car at 12 m/s; observer B is stationary. Consequently, sound reaching B from the rear and front speakers propagates right and left respectively. No heard frequency or beat-frequency value occurs in these data premises
\end{definition}
\begin{proof}
  This is the declared model definition of Matches Problem And Primary Figure.
\end{proof}
\begin{definition}[Uses Standard Still Air]
  \label{def:physics:phyx-mini-0222:phyxminiproblems-problemphyxmini0222-usesstandardstillair}
  \lean{PhyXMiniProblems.ProblemPhyXMini0222.UsesStandardStillAir}
  \uses{def:physics:phyx-mini-0222:phyxminiproblems-problemphyxmini0222-axialvelocityinmeterspersecond, def:physics:phyx-mini-0222:phyxminiproblems-problemphyxmini0222-speedinmeterspersecond, def:physics:phyx-mini-0222:phyxminiproblems-problemphyxmini0222-movingrailroadcarspeakersetup}
  The source gives no numerical sound speed and does not explicitly state the air motion. This separate standard-ambient calibration supplies still air and 343 m/s, rather than presenting either as a figure readout or as part of the requested conclusion
\end{definition}
\begin{proof}
  This is the declared model definition of Uses Standard Still Air.
\end{proof}
\begin{definition}[Has Physical Doppler Parameters]
  \label{def:physics:phyx-mini-0222:phyxminiproblems-problemphyxmini0222-hasphysicaldopplerparameters}
  \lean{PhyXMiniProblems.ProblemPhyXMini0222.HasPhysicalDopplerParameters}
  \uses{def:physics:phyx-mini-0222:phyxminiproblems-problemphyxmini0222-frequencyinhertz, def:physics:phyx-mini-0222:phyxminiproblems-problemphyxmini0222-speedinmeterspersecond, def:physics:phyx-mini-0222:phyxminiproblems-problemphyxmini0222-movingrailroadcarspeakersetup, def:physics:phyx-mini-0222:phyxminiproblems-problemphyxmini0222-speakervelocityrelativetoairinmeterspersecond}
  Positivity and subsonic conditions for the classical acoustic model
\end{definition}
\begin{proof}
  This is the declared model definition of Has Physical Doppler Parameters.
\end{proof}
\begin{definition}[Satisfies Directional Moving Source Doppler Laws]
  \label{def:physics:phyx-mini-0222:phyxminiproblems-problemphyxmini0222-satisfiesdirectionalmovingsourcedopplerlaws}
  \lean{PhyXMiniProblems.ProblemPhyXMini0222.SatisfiesDirectionalMovingSourceDopplerLaws}
  \uses{def:physics:phyx-mini-0222:phyxminiproblems-problemphyxmini0222-frequencyinhertz, def:physics:phyx-mini-0222:phyxminiproblems-problemphyxmini0222-speedinmeterspersecond, def:physics:phyx-mini-0222:phyxminiproblems-problemphyxmini0222-movingrailroadcarspeakersetup, def:physics:phyx-mini-0222:phyxminiproblems-problemphyxmini0222-speakervelocityrelativetoairinmeterspersecond, def:physics:phyx-mini-0222:phyxminiproblems-problemphyxmini0222-observerbvelocityrelativetoairinmeterspersecond}
  The one-dimensional classical Doppler laws for the two paths to stationary observer B. For the rear speaker the sound ray points right, giving f\_rear = f₀ (c - u\_B) / (c - u\_rear). For the front speaker the sound ray points left, so the signed velocity projections reverse and f\_front = f₀ (c + u\_B) / (c + u\_front). All velocities are relative to the air. These governing relations contain neither the requested numerical beat frequency nor an answer label
\end{definition}
\begin{proof}
  This is the declared model definition of Satisfies Directional Moving Source Doppler Laws.
\end{proof}
\begin{definition}[Satisfies Beat Frequency Law]
  \label{def:physics:phyx-mini-0222:phyxminiproblems-problemphyxmini0222-satisfiesbeatfrequencylaw}
  \lean{PhyXMiniProblems.ProblemPhyXMini0222.SatisfiesBeatFrequencyLaw}
  \uses{def:physics:phyx-mini-0222:phyxminiproblems-problemphyxmini0222-frequencyinhertz, def:physics:phyx-mini-0222:phyxminiproblems-problemphyxmini0222-movingrailroadcarspeakersetup}
  The acoustic definition of beat frequency: it is the magnitude of the difference between the two frequencies received by B. This law relates three independent physical frequency quantities without fixing a numerical answer
\end{definition}
\begin{proof}
  This is the declared model definition of Satisfies Beat Frequency Law.
\end{proof}
\begin{lemma}[heard Frequencies In Hertz eq exact Model Values]
  \label{lem:physics:phyx-mini-0222:phyxminiproblems-problemphyxmini0222-heardfrequenciesinhertz-eq-exactmodelvalues}
  \lean{PhyXMiniProblems.ProblemPhyXMini0222.heardFrequenciesInHertz_eq_exactModelValues}
  \uses{def:physics:phyx-mini-0222:phyxminiproblems-problemphyxmini0222-frequencyinhertz, def:physics:phyx-mini-0222:phyxminiproblems-problemphyxmini0222-movingrailroadcarspeakersetup, def:physics:phyx-mini-0222:phyxminiproblems-problemphyxmini0222-matchesproblemandprimaryfigure, def:physics:phyx-mini-0222:phyxminiproblems-problemphyxmini0222-usesstandardstillair, def:physics:phyx-mini-0222:phyxminiproblems-problemphyxmini0222-satisfiesdirectionalmovingsourcedopplerlaws}
  The two directional Doppler laws give the exact received-frequency readouts for the rear and front speakers under the explicit ambient calibration
\end{lemma}
\begin{proof}
  The conclusion follows from the governing relations and figure readouts named in the statement of heard Frequencies In Hertz eq exact Model Values.
\end{proof}
\begin{lemma}[beat Frequency In Hertz eq exact Model Value]
  \label{lem:physics:phyx-mini-0222:phyxminiproblems-problemphyxmini0222-beatfrequencyinhertz-eq-exactmodelvalue}
  \lean{PhyXMiniProblems.ProblemPhyXMini0222.beatFrequencyInHertz_eq_exactModelValue}
  \uses{def:physics:phyx-mini-0222:phyxminiproblems-problemphyxmini0222-frequencyinhertz, def:physics:phyx-mini-0222:phyxminiproblems-problemphyxmini0222-movingrailroadcarspeakersetup, def:physics:phyx-mini-0222:phyxminiproblems-problemphyxmini0222-matchesproblemandprimaryfigure, def:physics:phyx-mini-0222:phyxminiproblems-problemphyxmini0222-usesstandardstillair, def:physics:phyx-mini-0222:phyxminiproblems-problemphyxmini0222-satisfiesdirectionalmovingsourcedopplerlaws, def:physics:phyx-mini-0222:phyxminiproblems-problemphyxmini0222-satisfiesbeatfrequencylaw}
  Combining the two received frequencies with the beat-frequency law gives the exact classical-model beat readout, about 24.38 Hz
\end{lemma}
\begin{proof}
  The conclusion follows from the governing relations and figure readouts named in the statement of beat Frequency In Hertz eq exact Model Value.
\end{proof}
\begin{definition}[Answer Choice]
  \label{def:physics:phyx-mini-0222:phyxminiproblems-problemphyxmini0222-answerchoice}
  \lean{PhyXMiniProblems.ProblemPhyXMini0222.AnswerChoice}
  Labels of the four answers printed with the problem
\end{definition}
\begin{proof}
  This is the declared model definition of Answer Choice.
\end{proof}
\begin{definition}[hertz]
  \label{def:physics:phyx-mini-0222:phyxminiproblems-problemphyxmini0222-answerchoice-hertz}
  \lean{PhyXMiniProblems.ProblemPhyXMini0222.AnswerChoice.hertz}
  \uses{def:physics:phyx-mini-0222:phyxminiproblems-problemphyxmini0222-answerchoice}
  Whole-hertz beat frequency printed beside each displayed answer label
\end{definition}
\begin{proof}
  This is the declared model definition of hertz.
\end{proof}
\begin{definition}[recorded Answer Choice]
  \label{def:physics:phyx-mini-0222:phyxminiproblems-problemphyxmini0222-recordedanswerchoice}
  \lean{PhyXMiniProblems.ProblemPhyXMini0222.recordedAnswerChoice}
  \uses{def:physics:phyx-mini-0222:phyxminiproblems-problemphyxmini0222-answerchoice}
  The answer label recorded by the supplied dataset
\end{definition}
\begin{proof}
  This is the declared model definition of recorded Answer Choice.
\end{proof}
\begin{definition}[Matches Answer Choice]
  \label{def:physics:phyx-mini-0222:phyxminiproblems-problemphyxmini0222-matchesanswerchoice}
  \lean{PhyXMiniProblems.ProblemPhyXMini0222.MatchesAnswerChoice}
  \uses{def:physics:phyx-mini-0222:phyxminiproblems-problemphyxmini0222-acousticfrequency, def:physics:phyx-mini-0222:phyxminiproblems-problemphyxmini0222-frequencyinhertz, def:physics:phyx-mini-0222:phyxminiproblems-problemphyxmini0222-answerchoice}
  A physical beat frequency agrees with a displayed whole-hertz choice when its SI readout rounds to that integer. The tolerance is half a hertz
\end{definition}
\begin{proof}
  This is the declared model definition of Matches Answer Choice.
\end{proof}
% --- Archon declaration topology end ---
% --- Archon physics formalization source end ---
